\documentclass[10pt,a4paper]{amsart}
\usepackage[margin=19mm]{geometry}
\usepackage{amsmath,amssymb,mathtools}
\usepackage[scr=rsfs,cal=euler]{mathalfa}
\usepackage{xfrac}
\usepackage[unicode,hidelinks,bookmarks=false]{hyperref}
\newcommand{\ie}{i.e.\ }
\newcommand{\cf}{cf.\ }
\newcommand{\resp}{respectively\ }
\newcommand{\Subsection}{\subsection}
\providecommand{\nobreakdash}{\nobreak\hbox{-}}
\setlength{\emergencystretch}{2em}
\multlinegap0pt
\usepackage{amssymb}
\newtheorem{assumption}{Assumption}
\newtheorem{lemma}{Lemma}
\newtheorem{theorem}{Theorem}
\title{Rates of Convergence in the Central Limit Theorem for Markov Chains, with an Application to TD Learning}
\author{R. Srikant}
\date{Source: arXiv:2401.15719v5 (CC BY 4.0)}
\begin{document}
\maketitle

\noindent\emph{Source and licence.} Rates of Convergence in the Central Limit Theorem for Markov Chains, with an Application to TD Learning. Version arXiv:2401.15719v5 (CC BY 4.0), distributed by arXiv under the Creative Commons Attribution 4.0 International licence, http://creativecommons.org/licenses/by/4.0/, as recorded in the arXiv licence metadata for this exact version and shown on the abstract page https://arxiv.org/abs/2401.15719v5. Full licence notice: \texttt{licenses/arxiv-2401.15719-CC-BY-4.0.txt}.\par\smallskip

\noindent\textit{Editorial scope.} Counted unit: the article's theorem CLT (source0.tex lines 199--281). The article's complete original proof of it is delivered with it (source0.tex lines 283--747, one \texttt{proof} environment of 465 lines). No mathematics is written, repaired or re-derived: the statement and the proof are the article's own lines, in the article's own notation, and no retained line is substituted or re-flowed.  Retained uncounted as the source-local context the proof needs: lines 142--149; lines 154--156; lines 165--179.  Nothing was omitted from any retained span, so the package carries no omission record.  No retained line contains a citation, so the wrapper carries no bibliography and no bibliography entry.  No retained line contains a figure environment, an image-inclusion command or drawing code, so no image is delivered and the package declares no asset dependency.  Editorial formatting only, in addition: an AMS \texttt{amsart} preamble that restates the article's own declarations and macros used by the retained lines, namely the environments \texttt{assumption}, \texttt{lemma}, \texttt{theorem} on separate counters, together with the standard packages those lines need.  The retained environments print in this document as Assumption 1, Lemma 1, Theorem 1 in order of appearance, whereas the article prints the same items with the numbering of its own full document.  The complete native source of the article as distributed in the arXiv e-print for this exact version is bundled unchanged as \texttt{sources/153637/source0.tex}.
\begin{assumption}\label{assume: martingale}
Fix $\beta\in(0,1)$. We assume that the martingale difference sequence satisfies the following properties:
\begin{enumerate}\label{assume}
\item $E\|m_k\|^{2+\beta}<\infty$ for every $k\geq 1$.
\item $\Sigma_k=E(m_km_k^T\mid\mathcal{F}_{k-1})$ exists for every $k\geq 1$.
\end{enumerate}
\hfill $\diamond$
\end{assumption}

\begin{equation}\label{eq: normalized sum}
W_n:=\frac{S_n}{\sqrt{n}}
\end{equation}

\begin{lemma}\label{lem: regularity}
Let $g\in Lip_L$ and $f$ be the solution to Stein's equation:
\begin{equation}\label{eq: Stein}
    g(u)-E(g(Z))=\Delta f(u)-u^T\nabla f(u),
\end{equation}
where $\Delta$ and $\nabla$ are the Laplacian and gradient operators, respectively, and $Z\sim N(0,I)$. Then, $f$ has the following regularity properties for any $\beta\in (0,1)$:
\begin{enumerate}
    \item $||\nabla^2f(x)-\nabla^2f(y)||_{op}\leq \tilde{C}_1(d,\beta)||x-y||^\beta L,$ where $\nabla^2$ is the Hessian operator.
    \item $|\Delta f(x)-\Delta f(y)|\leq \tilde{C_2}(d,\beta)||x-y||^\beta L$,
where  $$\tilde{C}_1=C_1(d)+\frac{2}{1-\beta}, \quad \tilde{C}_2=C_2(d)+\frac{2d}{1-\beta}, \quad 
C_1=2^{3/2}\frac{(1+2d)\Gamma((1+d)/2)}{d\Gamma(d/2)},\quad C_2=2\sqrt{\frac{2d}{\pi}}.$$
    \item  $$\left|\langle\nabla^2 f(x),u_1u_2^T\rangle_{HS}\right|\leq C L||u_1||\cdot||u_2||\,\forall x, u_1, u_2\in \Re^d,$$ where $C$ is a constant (independent of $d$ and $\beta$) and $\langle\cdot,\cdot\rangle_{HS}$ is the Hilbert-Schmidt inner product (also called the Frobenius inner product), i.e., the sum of the element-by-element product of two matrices.
\end{enumerate}
\hfill $\diamond$
\end{lemma}

\begin{theorem}\label{thm: martingale CLT}
Let
\[
    V_k:=\sum_{j=k}^n E(\Sigma_j),
    \qquad k=1,\ldots,n,
\]
and assume that \(V_k\) is positive definite for each \(k=1,\ldots,n\).  Define
\[
    D_k:=V_k^{-1/2}E(\Sigma_k)V_k^{-1/2},
    \qquad k=1,\ldots,n .
\]
Let \(Z\sim N(0,I)\), and let \(W_n\) be defined as in
\((\ref{eq: normalized sum})\). Then, under
Assumption~\ref{assume: martingale} for the same fixed \(\beta\in(0,1)\),
\begin{align*}
   d_{\mathcal W}(W_n,n^{-1/2}V_1^{1/2}Z)
   &\le
   \frac{1}{\sqrt n}\sum_{k=1}^n
   \Biggl[
   \tilde C_1(d,\beta)\|V_k^{1/2}\|_{op}
   E\left(\|V_k^{-1/2}m_k\|^{2+\beta}\right)
   \\
   &\qquad\qquad
   +
   \tilde C_1(d,\beta)\operatorname{Tr}(D_k)\|V_k^{1/2}\|_{op}
   E\left(\|V_k^{-1/2}m_k\|^\beta\right)
   \Biggr]
   +\Delta_n,
\end{align*}
where
\begin{align*}
\Delta_n
:=
\sup_{h\in Lip_1}
\frac{1}{\sqrt n}
\left|
\sum_{k=1}^n
E\left[
\operatorname{Tr}\left(
A_{k,h}(\mathcal F_{k-1})
V_k^{-1/2}
\left(E(\Sigma_k)-\Sigma_k\right)
V_k^{-1/2}
\right)
\right]
\right|.
\end{align*}
Here
\[
A_{k,h}(\mathcal F_{k-1})
:=
E\left[
\nabla^2 f_{k,h}(\widetilde Z_k)
\mid
\mathcal F_{k-1}
\right],
\]
where \(f_{k,h}\) is the $\mathcal F_{k-1}$-measurable random solution of
the Stein equation defined in the proof below with
\[
    a=S_{k-1},
    \qquad
    Q_k=V_k^{1/2}.
\]
and
\[
    \widetilde Z_k
    :=
    V_k^{-1/2}T_{k+1}
    \sim
    N\left(0,V_k^{-1/2}V_{k+1}V_k^{-1/2}\right),
\]
with the convention \(V_{n+1}=0\). Moreover,
\[
    \|A_{k,h}(\mathcal F_{k-1})\|_{op}
    \le
    C\|V_k^{1/2}\|_{op}, \quad \text{a.s.},
\]
where \(C\) is the constant in statement (3) of Lemma~\ref{lem: regularity}.
Since \(0\preceq D_k\preceq I\), the same bound also implies the coarser
version obtained by replacing \(\operatorname{Tr}(D_k)\) by \(d\).  The sharper
form above is used in the Markov-chain applications.
\end{theorem}

\begin{proof}
Fix \(h\in Lip_1\). By the standard rescaling of Lipschitz test functions,
\[
d_{\mathcal W}(W_n,n^{-1/2}V_1^{1/2}Z)
=
\sup_{h\in Lip_1}
\frac{1}{\sqrt n}
E\left[h(S_n)-h(T_1)\right],
\]
where
\[
    T_k:=\sum_{j=k}^n Y_j,
\]
and \(Y_1,\ldots,Y_n\) are independent Gaussian random vectors, independent
of the martingale difference sequence, with
\[
    Y_j\sim N(0,E(\Sigma_j)).
\]
Thus
\[
    T_k\sim N(0,V_k).
\]
The rescaling follows because if \(g\in Lip_1\), then
\(h(x):=\sqrt n\,g(x/\sqrt n)\) also belongs to \(Lip_1\), and
\[
E\left[g(W_n)-g(n^{-1/2}V_1^{1/2}Z)\right]
=
\frac{1}{\sqrt n}
E\left[h(S_n)-h(T_1)\right].
\]

Using the Lindeberg decomposition,
\[
E\left[h(S_n)-h(T_1)\right]
=
\sum_{k=1}^n
E\left[
h(S_k+T_{k+1})-h(S_{k-1}+T_k)
\right],
\]
where \(T_{n+1}=S_0=0\).

We now bound a single summand. Fix \(k\). Conditional on
\(\mathcal F_{k-1}\), set
\[
    a=S_{k-1},
    \qquad
    Q_k=V_k^{1/2}.
\]
Define
\[
    \widetilde h_k(u):=h(a+Q_k u).
\]
Since \(h\in Lip_1\), we have
\[
    \widetilde h_k\in Lip_{L_k},
    \qquad
    L_k=\|Q_k\|_{op}=\|V_k^{1/2}\|_{op}.
\]
Let \(f_k=f_{k,h}\) be the solution of the Stein equation
\[
    \widetilde h_k(u)-E[\widetilde h_k(Z)]
    =
    \Delta f_k(u)-u^\top\nabla f_k(u).
\]
Since \(a=S_{k-1}\) is \(\mathcal F_{k-1}\)-measurable, the function
\(f_k\) is also \(\mathcal F_{k-1}\)-measurable as a random function.

By the Stein equation, with
\[
    \widetilde Z_k:=V_k^{-1/2}T_{k+1},
\]
we have
\begin{align*}
& E\left[
h(S_k+T_{k+1})-h(S_{k-1}+T_k)
\mid
\mathcal F_{k-1}
\right]
\\
&=
E\biggl[
\Delta f_k\left(V_k^{-1/2}m_k+\widetilde Z_k\right)
-
\left(V_k^{-1/2}m_k+\widetilde Z_k\right)^\top
\nabla f_k\left(V_k^{-1/2}m_k+\widetilde Z_k\right)
\mid
\mathcal F_{k-1}
\biggr].
\end{align*}
Moreover,
\[
    \widetilde Z_k
    \sim
    N(0,\Lambda_k),
    \qquad
    \Lambda_k:=V_k^{-1/2}V_{k+1}V_k^{-1/2}.
\]
Since
\[
    V_k=E(\Sigma_k)+V_{k+1},
\]
we have
\[
    I-\Lambda_k
    =
    V_k^{-1/2}E(\Sigma_k)V_k^{-1/2}.
\]

We rewrite the preceding conditional expectation as
\begin{align*}
& E\biggl[
\operatorname{Tr}\left(
\Lambda_k
\nabla^2 f_k\left(V_k^{-1/2}m_k+\widetilde Z_k\right)
\right)
-
\widetilde Z_k^\top
\nabla f_k\left(V_k^{-1/2}m_k+\widetilde Z_k\right)
\mid
\mathcal F_{k-1}
\biggr]
\\
&+
E\biggl[
\operatorname{Tr}\left(
(I-\Lambda_k)
\nabla^2 f_k\left(V_k^{-1/2}m_k+\widetilde Z_k\right)
\right)
\mid
\mathcal F_{k-1}
\biggr]
\\
&-
E\left[
(V_k^{-1/2}m_k)^\top\nabla f_k(\widetilde Z_k)
\mid
\mathcal F_{k-1}
\right]
\\
&-
E\biggl[
(V_k^{-1/2}m_k)^\top
\left(
\nabla f_k\left(V_k^{-1/2}m_k+\widetilde Z_k\right)
-
\nabla f_k(\widetilde Z_k)
\right)
\mid
\mathcal F_{k-1}
\biggr].
\end{align*}
The first term is zero by Gaussian integration by parts: conditional on
$\mathcal F_k$, the quantity $V_k^{-1/2}m_k$ is fixed and
$\widetilde Z_k\sim N(0,\Lambda_k)$ is independent of $\mathcal F_k$. The third term is zero because $f_k$ is $\mathcal F_{k-1}$-measurable as a
random function, $\widetilde Z_k$ is independent of $\mathcal F_k$, and
$E[m_k\mid\mathcal F_{k-1}]=0$.

Using the fundamental theorem of calculus, we obtain
\begin{align*}
& E\left[
h(S_k+T_{k+1})-h(S_{k-1}+T_k)
\mid
\mathcal F_{k-1}
\right]
\\
&=
E\biggl[
\operatorname{Tr}\left(
V_k^{-1/2}E(\Sigma_k)V_k^{-1/2}
\nabla^2 f_k\left(V_k^{-1/2}m_k+\widetilde Z_k\right)
\right)
\mid
\mathcal F_{k-1}
\biggr]
\\
&\quad
-
E\biggl[
(V_k^{-1/2}m_k)^\top
\left(
\nabla^2 f_k(\Theta V_k^{-1/2}m_k+\widetilde Z_k)
-
\nabla^2 f_k(\widetilde Z_k)
\right)
(V_k^{-1/2}m_k)
\mid
\mathcal F_{k-1}
\biggr]
\\
&\quad
-
E\biggl[
(V_k^{-1/2}m_k)^\top
\nabla^2 f_k(\widetilde Z_k)
(V_k^{-1/2}m_k)
\mid
\mathcal F_{k-1}
\biggr],
\end{align*}
where \(\Theta\sim {\rm Uniform}[0,1]\) is independent of all other random
variables.

Define
\[
    A_{k,h}(\mathcal F_{k-1})
    :=
    E\left[
    \nabla^2 f_k(\widetilde Z_k)
    \mid
    \mathcal F_{k-1}
    \right].
\]
Using
\[
    \Sigma_k=E[m_km_k^\top\mid\mathcal F_{k-1}],
\]
the last term becomes
\[
    -
    \operatorname{Tr}\left(
    V_k^{-1/2}\Sigma_k V_k^{-1/2}
    A_{k,h}(\mathcal F_{k-1})
    \right).
\]
By statement (1) of Lemma~\ref{lem: regularity}, the second term above is
bounded in absolute value by
\[
\tilde C_1(d,\beta)\|V_k^{1/2}\|_{op}
E\left[
\|V_k^{-1/2}m_k\|^{2+\beta}
\mid
\mathcal F_{k-1}
\right].
\]
Therefore,
\begin{align*}
& E\left[
h(S_k+T_{k+1})-h(S_{k-1}+T_k)
\mid
\mathcal F_{k-1}
\right]
\\
&\le
E\biggl[
\operatorname{Tr}\left(
V_k^{-1/2}E(\Sigma_k)V_k^{-1/2}
\nabla^2 f_k\left(V_k^{-1/2}m_k+\widetilde Z_k\right)
\right)
\mid
\mathcal F_{k-1}
\biggr]
\\
&\quad
-
\operatorname{Tr}\left(
V_k^{-1/2}\Sigma_k V_k^{-1/2}
A_{k,h}(\mathcal F_{k-1})
\right)
\\
&\quad
+
\tilde C_1(d,\beta)\|V_k^{1/2}\|_{op}
E\left[
\|V_k^{-1/2}m_k\|^{2+\beta}
\mid
\mathcal F_{k-1}
\right].
\end{align*}

Next, add and subtract
\[
\operatorname{Tr}\left(
V_k^{-1/2}E(\Sigma_k)V_k^{-1/2}
A_{k,h}(\mathcal F_{k-1})
\right).
\]
This gives
\begin{align*}
& E\left[
h(S_k+T_{k+1})-h(S_{k-1}+T_k)
\mid
\mathcal F_{k-1}
\right]
\\
&\le
\operatorname{Tr}\left(
A_{k,h}(\mathcal F_{k-1})
V_k^{-1/2}
\left(E(\Sigma_k)-\Sigma_k\right)
V_k^{-1/2}
\right)
\\
&\quad
+
E\biggl[
\operatorname{Tr}\left(
D_k
\left[
\nabla^2 f_k\left(V_k^{-1/2}m_k+\widetilde Z_k\right)
-
\nabla^2 f_k(\widetilde Z_k)
\right]
\right)
\mid
\mathcal F_{k-1}
\biggr]
\\
&\quad
+
\tilde C_1(d,\beta)\|V_k^{1/2}\|_{op}
E\left[
\|V_k^{-1/2}m_k\|^{2+\beta}
\mid
\mathcal F_{k-1}
\right],
\end{align*}
where
\[
    D_k:=V_k^{-1/2}E(\Sigma_k)V_k^{-1/2}.
\]

It remains to bound the weighted trace term. Since
\[
    V_k=E(\Sigma_k)+V_{k+1},
\]
we have
\[
    D_k=I-V_k^{-1/2}V_{k+1}V_k^{-1/2}.
\]
Hence
\[
    0\preceq D_k\preceq I.
\]
Therefore, for any \(x,u\in\mathbb R^d\),
\begin{align*}
&\left|
\operatorname{Tr}\left(
D_k
\left[
\nabla^2 f_k(x+u)-\nabla^2 f_k(x)
\right]
\right)
\right|
\\
&\le
\operatorname{Tr}(D_k)
\left\|
\nabla^2 f_k(x+u)-\nabla^2 f_k(x)
\right\|_{op}
\\
&\le
   \operatorname{Tr}(D_k)\,\tilde C_1(d,\beta)\|V_k^{1/2}\|_{op}\|u\|^\beta.
\end{align*}
Applying this bound with
\[
    x=\widetilde Z_k,
    \qquad
    u=V_k^{-1/2}m_k,
\]
and taking conditional expectation gives
\begin{align*}
&\left|
E\biggl[
\operatorname{Tr}\left(
D_k
\left[
\nabla^2 f_k\left(V_k^{-1/2}m_k+\widetilde Z_k\right)
-
\nabla^2 f_k(\widetilde Z_k)
\right]
\right)
\mid
\mathcal F_{k-1}
\biggr]
\right|
\\
&\le
\tilde C_1(d,\beta)\operatorname{Tr}(D_k)\|V_k^{1/2}\|_{op}
E\left[
\|V_k^{-1/2}m_k\|^\beta
\mid
\mathcal F_{k-1}
\right].
\end{align*}

Thus,
\begin{align*}
& E\left[
h(S_k+T_{k+1})-h(S_{k-1}+T_k)
\mid
\mathcal F_{k-1}
\right]
\\
&\le
\operatorname{Tr}\left(
A_{k,h}(\mathcal F_{k-1})
V_k^{-1/2}
\left(E(\Sigma_k)-\Sigma_k\right)
V_k^{-1/2}
\right)
\\
&\quad
+
\tilde C_1(d,\beta)\operatorname{Tr}(D_k)\|V_k^{1/2}\|_{op}
E\left[
\|V_k^{-1/2}m_k\|^\beta
\mid
\mathcal F_{k-1}
\right]
\\
&\quad
+
\tilde C_1(d,\beta)\|V_k^{1/2}\|_{op}
E\left[
\|V_k^{-1/2}m_k\|^{2+\beta}
\mid
\mathcal F_{k-1}
\right].
\end{align*}

Taking expectations, summing over \(k\), substituting into the Lindeberg
decomposition, and using the rescaling identity for the Wasserstein distance,
we obtain, for the fixed \(h\in Lip_1\),
\begin{align*}
& E\left[
h(W_n)-h(n^{-1/2}V_1^{1/2}Z)
\right]
\\
&\le
\frac{1}{\sqrt n}\sum_{k=1}^n
\Biggl[
\tilde C_1(d,\beta)\|V_k^{1/2}\|_{op}
E\left(\|V_k^{-1/2}m_k\|^{2+\beta}\right)
\\
&\qquad\qquad
+
\tilde C_1(d,\beta)\operatorname{Tr}(D_k)\|V_k^{1/2}\|_{op}
E\left(\|V_k^{-1/2}m_k\|^\beta\right)
\Biggr]
\\
&\quad
+
\frac{1}{\sqrt n}
\left|\sum_{k=1}^n
E\left[
\operatorname{Tr}\left(
A_{k,h}(\mathcal F_{k-1})
V_k^{-1/2}
\left(E(\Sigma_k)-\Sigma_k\right)
V_k^{-1/2}
\right)
\right]
\right|.
\end{align*}
Taking the supremum over \(h\in Lip_1\) proves the claimed bound.

Finally, statement (3) of Lemma~\ref{lem: regularity} gives
\[
    \|A_{k,h}(\mathcal F_{k-1})\|_{op}
    \le
    C\|V_k^{1/2}\|_{op}.
\]
This completes the proof.
\end{proof}

\end{document}
